% /root/monodromy, 19 November 2026; independently audited by /root.
\subsection{A terminating finite-field algorithm for the entire monodromy image}
\label{subsec:effective-entire-monodromy}
\status{Proved; independent parent proof audit and exact quartic computation passed.
This is an effective answer for specified finite-field inputs, not a structural
classification for arbitrary perfect base fields.}

\begin{theorem}\label{thm:effective-monodromy-matrices}
Suppose $k=\mathbf F_q$ and the fixed smooth projective
embedding $X\subset\mathbf P^m_k$ is given by homogeneous
equations.  For every
\[
 d\geq d_{\rm coh}=\max\{1,\operatorname{reg}(\omega_X)-1\},
\]
there is an explicit terminating algorithm that outputs,
up to a simultaneous choice of basis, the complete finite
lists of matrices
\[
 G_{\rm geom},G_{\rm arith}\subset\GL_{r_d}(\mathbf F_p)
\]
for the local system in A.4.  It also decides geometric
and arithmetic irreducibility and outputs the determinant
images.  The algorithm uses only finite linear algebra,
Groebner-basis computations, construction of finite fields,
and bounded enumeration of polynomials and matrices.
It does not call an oracle for Galois groups or for
factorization over an algebraic closure.
\end{theorem}

\begin{proof}
We give the steps and their termination proofs.

\emph{Step 1: the stable matrix.}
Use Proposition~\ref{prop:effective-generic-rank} over
the rational function field $K=\mathbf F_q(t_1,\ldots,t_M)$
of an affine chart of the ambient degree-$d$ parameter
space.  Let $B$ be its coherent Frobenius matrix and
$b$ the coherent dimension.  Compute the image of
$B B^{[p]}\cdots B^{[p^{b-1}]}$ by Gaussian elimination.
The Frobenius-linearized map on that image is invertible;
a chosen basis gives an explicit $A\in\GL_r(K)$, where
$r=r_d$.  If $r=0$, output the trivial groups and the zero
representation (which is not called irreducible), and stop.
This construction is valid even though $K$ is imperfect.

\emph{Step 2: the frame algebra.}
Let $T=(T_{ij})$ be an $r\times r$ matrix of indeterminates
and introduce one more indeterminate $z$.  Form
\begin{equation}\label{eq:effective-frame-algebra}
 R=K[T_{ij},z]/
        (A T^{[p]}-T,\ z\det T-1),\qquad
 D=\prod_{i=0}^{r-1}(p^r-p^i)=|\GL_r(\mathbf F_p)|.
\end{equation}
This is a finite etale $K$-algebra of dimension $D$.
Indeed the equations for one column, $Av^{[p]}-v=0$,
have Jacobian $-1$.  Their leading homogeneous forms
are an invertible linear combination of the $p$-th
powers of the coordinates, so they have no common
projective zero at infinity and define a scheme of
length $p^r$.  Its geometric points form an
$\mathbf F_p$-space of dimension $r$ and span the
whole $r$-dimensional vector space: otherwise their
span, of dimension $s<r$, would itself have only
$p^s$ fixed points by the same argument.
The invertible frames are thus a torsor under
$\GL_r(\mathbf F_p)$ and number exactly $D$.
Compute a Groebner basis of the displayed ideal,
a $K$-basis of $R$, and its multiplication table.

\emph{Step 3: a primitive element from a finite list.}
Let $e_1,\ldots,e_D$ be the computed basis and put
$N=D(D-1)/2$.  Choose the finite set
$S=\{1,t_1,t_1^2,\ldots,t_1^N\}\subset K$.
Test the elements
\[
 \theta=\sum_{j=1}^D a_j e_j,\qquad (a_1,\ldots,a_D)\in S^D,
\]
until the vectors $1,\theta,\ldots,\theta^{D-1}$ are
linearly independent.  At least one test succeeds.
Over a separable closure, the algebra has $D$ distinct
embeddings.  The product of the pairwise differences
of their values at $\sum a_je_j$ is a nonzero polynomial
of degree $N$ in the $a_j$.  A nonzero polynomial of
degree at most $N$ cannot vanish on a grid having more
than $N$ distinct values in every coordinate, by
induction on the number of variables.  The successful
$\theta$ has a square-free monic minimal polynomial
$f\in K[Z]$ of degree $D$, computed by linear algebra,
and $R=K[\theta]\simeq K[Z]/(f)$.

\emph{Step 4: factorization by bounded enumeration.}
Clear the denominators of $f$ to obtain a nonzero
$F\in\mathbf F_q[t_1,\ldots,t_M,Z]$.
Enumerate all normalized polynomials over $\mathbf F_q$
of total degree at most $\deg F$, and test divisibility
by polynomial division or Groebner reduction.
This is a finite list.  The irreducible factors of
$F$ are among it, since total degrees add in a
polynomial domain.  Repeated exact division therefore
factors $F$ completely.  Discard factors independent
of $Z$, which are units over $K$, and normalize the
remaining ones to be monic in $Z$.
Gauss's lemma proves that these are precisely the
irreducible factors of $f$ over $K$.
Thus this step is a fully specified finite algorithm;
practical multivariate factorization can replace the
enumeration without changing its correctness.

Choose one factor $f_0$ by a fixed lexicographic rule.
It defines a field component $R_0$ of $R$ and the
corresponding maximal ideal
\[
 I_0=(A T^{[p]}-T,\ z\det T-1,\ f_0(\theta))
       \subset K[T_{ij},z].
\]

\emph{Step 5: the arithmetic matrices.}
Enumerate all $g\in\GL_r(\mathbf F_p)$.
Its action on the frame algebra is the substitution
\[
 T\longmapsto Tg,\qquad z\longmapsto(\det g)^{-1}z.
\]
Retain exactly those $g$ for which $f_0(\theta)$ after
this substitution reduces to zero modulo $I_0$.
This is a finite collection of Groebner reductions.
The retained list is the stabilizer of the chosen
field component, hence exactly $G_{\rm arith}$ in
the corresponding frame.  Preservation of the one
extra equation suffices because the frame ideal is
already preserved; inclusion between the resulting
maximal ideals is equality.

To justify the monodromy identification explicitly,
choose a geometric frame in the chosen component.
Its Galois orbit is the image of $\operatorname{Gal}(K)$
acting on the frames.  Under the simply transitive
right frame action, the stabilizer of its connected
component is the same subgroup of $\GL_r(\mathbf F_p)$,
up to the common initial choice of frame.
Restricting the original local system to any dense
open does not change its image: the fundamental group
of a dense open in a normal connected scheme surjects
onto that of the scheme.  Thus the generic calculation
computes the image over the stated $U_d$.

\emph{Step 6: the geometric matrices with a finite
constant-field bound.}
Construct the finite field $k_D=\mathbf F_{q^D}$ and
put $K_D=k_D(t_1,\ldots,t_M)$.
A finite field of this degree is constructed by
enumerating monic degree-$D$ polynomials and testing
irreducibility; the elementary existence theorem
for finite fields guarantees termination.
Factor $f_0$ over $K_D$ by Step 4, now over $k_D$,
choose one factor $f_{00}$, and repeat the matrix
preservation test with its component ideal.
The resulting subgroup is $G_{\rm geom}$, and it is
a subgroup of the already chosen arithmetic image.

Here is the reason the finite extension degree $D$
suffices.  The selected field component is a connected
torsor under its stabilizer $H$, so its degree over
$K$ is $|H|$, a divisor of $D$.  Its full algebraic
constant field is some $\mathbf F_{q^c}$, where
$c\mid |H|$, hence $c\mid D$.
After extending to $k_D$, every resulting field
component has its constant field algebraically closed
in it.  A finitely generated field over a perfect
field is separably generated; therefore each such
component is geometrically integral.
Consequently these are exactly the geometric
components, and their stabilizers give the geometric
image, not an upper approximation to it.

\emph{Step 7: irreducibility and determinant.}
For each of the two finite matrix lists, enumerate
all row-echelon matrices representing the proper
nonzero subspaces of $\mathbf F_p^r$ and test
invariance under every listed element.
There is an invariant subspace if and only if the
representation is reducible.  Taking determinants
of the listed matrices gives both determinant
images.  Proposition~\ref{prop:exact-general-determinant-image}
is a much smaller independent calculation of the
same images and provides a consistency check.

Every step has a finite bound or an explicitly proved
termination criterion.  This proves the theorem.
\end{proof}

\begin{remark}\label{rem:effective-monodromy-scope}
The theorem outputs the whole finite group for each
specified finite-field input, with a proved finite
algorithm.  Its worst-case complexity is enormous:
even the frame dimension is $|\GL_r(\mathbf F_p)|$.
It does not supply a uniform structural description
of the groups as $X$ and $d$ vary.  Nor does it supply
a uniform algorithm for arithmetic monodromy over
an arbitrary perfect field presented without effective
field operations and factorization data.
Those distinctions are retained in the result ledger.
\end{remark}

\begin{example}\label{ex:effective-quartic-frame-algebra}
The retained program \path{lang_frame_quartic.py}
executes a nontrivial small frame calculation for a genuine
smooth plane-quartic family.  Over $\mathbf F_2(t)$ take
\[
 f_t=x^4+y^4+z^4+x^2yz+xz^3+txy^3+xyz^2.
\]
In the standard genus-three monomial basis its Hasse matrix is
\[
 A_t=\begin{pmatrix}1&0&0\\0&0&1\\0&t&1\end{pmatrix}.
\]
The program extracts these coefficients from $f_t$ and checks
that all three affine-chart singularity ideals have Groebner
basis $[1]$ over $\mathbf F_2(t)$.
Thus the generic curve is smooth and ordinary, since
$\det A_t=t\ne0$.

The first coordinate is a constant summand.  A frame of the
remaining rank-two Lang module has columns $(c^2,c)^t$ and
$(d^2,d)^t$, with algebra
\[
 R=\mathbf F_2(t)[c,d]/
 \left(c^3+\frac{c+1}{t},\quad
              d^2+cd+c^2+\frac1t\right).
\]
Its basis is $1,c,c^2,d,cd,c^2d$, and its frame determinant
is $cd(c+d)=1/t$.
The program verifies all six substitutions from
$\GL_2(\mathbf F_2)$ preserve this algebra.
It also finds the primitive element $\theta=c+td$,
whose minimal polynomial, after clearing denominators, is
\[
 t^2 Z^6+(t^3+t^2)Z^3+(t^4+t^2+1)Z^2
       +(t^4+t)Z+t^6+t^5+1.
\]
Linear independence of $1,\theta,\ldots,\theta^5$ and
annihilation by this polynomial are checked exactly.

Its geometric irreducibility has a direct certificate.
The first equation gives $t=(c+1)/c^3$, a separable
degree-three map of rational function fields.
Setting $u=d/c$ transforms the second equation into
\[
 u^2+u=\frac1{c+1}.
\]
The right side has a simple pole and is not an
Artin--Schreier coboundary over $\overline{\mathbf F}_2(c)$:
every pole of $v^2+v$ has even order.
Thus the frame algebra remains a field of degree six
after extending constants to the algebraic closure.
Both monodromy images of this particular curve family
are exactly
$\operatorname{diag}(1,\GL_2(\mathbf F_2))\subset\GL_3(\mathbf F_2)$.
This is a restricted one-parameter family, not the full
space of plane quartics, whose image was determined
earlier.  Its exact coefficient matrix, basis, six group
elements, minimal polynomial and smoothness certificates
are retained in \texttt{lang\_frame\_quartic.json}.
\end{example}
